% cellular-network-performance.tex — why requests are slow and expensive on cellular: RRC radio states
% (LTE idle/connected + DRX, 3G DCH/FACH, 5G RRC_INACTIVE), promotion delay and tail energy with
% measured numbers, latency by generation, the round trips of a cold HTTPS request (DNS, TCP, TLS,
% slow start) vs a reused connection vs HTTP/3, connection reuse + HTTP/2 coalescing, prefetch /
% batching / priorities, adapting (Low Data Mode, waitsForConnectivity, timeouts), measuring
% (URLSessionTaskMetrics phases, Instruments HTTP Traffic, Network Link Conditioner).
% Picture: LTE RRC state machine + the cold-vs-warm request timeline drawn to scale.
% Sources (checked 2026-09-25): hpbn.co/mobile-networks/ (Tables 7-2, 7-9; 3G promotions); Huang et al.,
% "A Close Examination of Performance and Power Characteristics of 4G LTE Networks", MobiSys 2012
% (feng-qian.github.io/paper/lte_mobisys12.pdf); 3GPP TS 36.331 / 38.331; RFC 9113 §9.1.1, RFC 9110
% §15.5.20 (421), RFC 6928; developer.apple.com (URLSessionTaskTransactionMetrics,
% URLSessionConfiguration, URLSessionTask.priority); developer.apple.com/videos/play/wwdc2021/10212;
% developer.apple.com/forums/thread/110187 (connection pool timeout is not API).
% NOT repeated: tail-energy batching diagram, isDiscretionary, MetricKit list
% (ios-platform/app-size-and-energy.tex); HTTP/2 vs /3 stacks, caching (networking/http-deep.tex).
% @labels: area=networking kind=concept platform=ios topic=networking,performance discussion=218
% @tags: cellular, rrc-states, tail-energy, promotion-delay, round-trips, tcp-slow-start, connection-reuse, http2-coalescing, prefetching, urlsessiontaskmetrics, network-link-conditioner, low-data-mode
\documentclass[8pt]{extarticle}
\usepackage{printup-sheet}
\usepackage{array}

\lstdefinelanguage{SwiftSheet}{
  morekeywords={protocol,class,final,struct,enum,func,var,let,weak,init,for,in,where,
    if,else,return,guard,self,nil,try,await,async,throws,private,some,
    true,false,Int,Date},
  sensitive=true, morecomment=[l]{//}, morestring=[b]",
  literate={->}{{\hbox{-}\hbox{>}}}2 {==}{{\hbox{=}\hbox{=}}}2 {??}{{\hbox{?}\hbox{?}}}2}
\lstset{basicstyle=\ttfamily\scriptsize, aboveskip=2pt, belowskip=2pt}
\newcommand\ct[1]{\texttt{#1}}
\newcommand\hd[1]{\par\vspace{3pt}\noindent{\bfseries\color{sheetBlue}#1}\par\vspace{1pt}}

\begin{document}

\sheettitle{Network performance on cellular — radio states, round trips, reuse}{networking · folio}

\oneliner{On cellular a small request is dominated not by bytes but by \textbf{waking the radio}
(promotion from idle: $\sim$0.1–0.3\,s on LTE, up to 2\,s on 3G, then a \textbf{tail} of seconds at high
power) and by \textbf{round trips}: DNS + TCP + TLS before the first byte, each one cellular
\textbf{RTT}. So: \textbf{fewer, earlier, bigger} requests over connections you \textbf{already have} —
and measure the phases, not the total.}

\vspace{3pt}
\noindent\begin{tikzpicture}[sheet,
    st/.style={box, font=\scriptsize, inner sep=2pt, align=center},
    l/.style={font=\tiny, text=black!75, inner sep=1pt, align=center},
    sg/.style={draw=white, fill=#1, minimum height=3.6mm, inner sep=0pt, font=\tiny, anchor=west}]
  % ---------- RRC state machine ----------
  \node[font=\bfseries\small, text=sheetBlue, anchor=west] at (-0.3,3.7) {LTE radio (RRC) states};
  \node[st, draw=sheetGrey, fill=black!4, minimum width=17mm, minimum height=9mm] (idle) at (0.7,1.6) {\textbf{RRC\_IDLE}\\\tiny paging only\\\tiny $\approx$ no power};
  \node[st, draw=sheetRed, fill=sheetRed!6, text width=31mm, minimum height=20mm] (con) at (4.65,1.6) {};
  \node[font=\scriptsize\bfseries, anchor=north] at (4.65,2.55) {RRC\_CONNECTED};
  \node[st, draw=sheetRed, fill=sheetRed!15, font=\tiny] (cr) at (4.65,1.95) {continuous reception (data)};
  \node[st, draw=sheetOrange, fill=sheetOrange!15, font=\tiny] (sd) at (4.65,1.35) {short DRX: sleep, wake 20 ms};
  \node[st, draw=sheetOrange, fill=sheetOrange!8, font=\tiny] (ld) at (4.65,0.8) {long DRX: wake every 40 ms};
  \draw[->, thin, sheetGrey] (cr) -- (sd); \draw[->, thin, sheetGrey] (sd) -- (ld);
  \draw[hot, draw=sheetRed] ([yshift=0.25cm]idle.east) -- node[l, above, text=sheetRed]{data $\to$\\promotion} node[l, below=2pt, text=sheetRed]{$\sim$260 ms} ([yshift=0.25cm]con.west);
  \draw[hot, draw=sheetGreen!70!black] ([yshift=-0.55cm]con.west) -- node[l, below, text=sheetGreen!50!black, align=center]{inactivity timer\\\textbf{tail $\approx$ 11.5 s}} ([yshift=-0.55cm]idle.east);
  \node[l, anchor=west, align=left] at (-0.3,0.05) {Measured on one LTE network (Huang et al., MobiSys 2012): promotion 260 ms, tail 11.576 s; one packet\\
     costs \textbf{12.8 J on LTE} vs 7.4 J on 3G vs 0.04 J on Wi-Fi. Timers are set by the network.};
  \node[l, anchor=west, align=left, text=sheetBrown] at (-0.3,-0.55) {3G: IDLE $\to$ DCH up to 2 s, FACH $\to$ DCH up to 1.5 s (HPBN). 5G NR adds \textbf{RRC\_INACTIVE}:\\context kept, faster resume.};
  \draw[sheetGrey!40] (6.95,3.8) -- (6.95,-0.8);
  % ---------- request timeline, 1 cm = 80 ms = 1 RTT ----------
  \node[font=\bfseries\small, text=sheetBlue, anchor=west] at (7.05,3.7) {One small GET, RTT 80 ms — drawn to scale};
  \foreach \x/\t in {8.4/0, 10.9/200, 13.4/400, 15.9/600} {\draw[sheetGrey!60] (\x,3.25) -- (\x,0.35); \node[l, above] at (\x,3.25) {\t\,ms};}
  % cold h2
  \node[l, anchor=east, align=right] at (8.35,2.9) {cold radio,\\TCP + TLS 1.3};
  \node[sg=sheetRed!45, minimum width=32.5mm] at (8.4,2.9) {promotion};
  \node[sg=sheetBrown!35, minimum width=10mm] at (11.65,2.9) {DNS};
  \node[sg=sheetBlue!30, minimum width=10mm] at (12.65,2.9) {TCP};
  \node[sg=sheetGreen!35, minimum width=10mm] at (13.65,2.9) {TLS};
  \node[sg=sheetOrange!45, minimum width=10mm] at (14.65,2.9) {GET};
  \node[sg=black!15, minimum width=6mm] at (15.65,2.9) {srv};
  \node[l, anchor=west, text=sheetRed] at (16.3,2.9) {$\approx$630};
  % cold h3
  \node[l, anchor=east, align=right] at (8.35,2.2) {cold radio,\\HTTP/3 (QUIC)};
  \node[sg=sheetRed!45, minimum width=32.5mm] at (8.4,2.2) {promotion};
  \node[sg=sheetBrown!35, minimum width=10mm] at (11.65,2.2) {DNS};
  \node[sg=sheetGreen!35, minimum width=10mm] at (12.65,2.2) {QUIC+TLS};
  \node[sg=sheetOrange!45, minimum width=10mm] at (13.65,2.2) {GET};
  \node[sg=black!15, minimum width=6mm] at (14.65,2.2) {srv};
  \node[l, anchor=west] at (15.3,2.2) {$\approx$550 (0-RTT: 470)};
  % warm radio, new connection
  \node[l, anchor=east, align=right] at (8.35,1.5) {radio up,\\new connection};
  \node[sg=sheetBrown!35, minimum width=10mm] at (8.4,1.5) {DNS};
  \node[sg=sheetBlue!30, minimum width=10mm] at (9.4,1.5) {TCP};
  \node[sg=sheetGreen!35, minimum width=10mm] at (10.4,1.5) {TLS};
  \node[sg=sheetOrange!45, minimum width=10mm] at (11.4,1.5) {GET};
  \node[sg=black!15, minimum width=6mm] at (12.4,1.5) {srv};
  \node[l, anchor=west] at (13.05,1.5) {$\approx$370};
  % warm reused
  \node[l, anchor=east, align=right] at (8.3,0.8) {radio up,\\\textbf{reused} conn.};
  \node[sg=sheetOrange!45, minimum width=10mm] at (8.4,0.8) {GET};
  \node[sg=black!15, minimum width=6mm] at (9.4,0.8) {srv};
  \node[l, anchor=west, text=sheetGreen!50!black] at (10.05,0.8) {\textbf{$\approx$130} — 5$\times$ faster: reuse is the win};
  \node[l, anchor=west, align=left] at (7.05,0.05) {then the radio stays up $\approx$11.5 s (tail) whatever you do — so \textbf{piggyback}:\\send queued analytics and prefetch \emph{now}, while the tail is paid for.};
  \node[l, anchor=west, align=left, text=sheetBrown] at (7.05,-0.55) {HPBN, active connection: 2G 300–1000 ms · 3G 100–500 ms · 4G $<$ 100 ms RTT.\\A 1 MB body adds several RTTs of TCP slow start (initial window $\approx$ 10 segments).};
\end{tikzpicture}

\begin{multicols}{2}
\raggedright\setstretch{1.0}

\hd{How it works}
\begin{itemize}
  \item \textbf{Radio states} (RRC): the phone holds a dedicated channel only in CONNECTED. Getting
        there is a signalling exchange (\textbf{promotion}); leaving is on an \textbf{inactivity
        timer} set by the carrier, and the radio burns near-peak power during that \textbf{tail}.
        Inside CONNECTED, DRX sleeps between short wake-ups.
  \item \textbf{Consequence}: a 1 KB ping every 30 s keeps the radio up $\sim$40\,\% of the time
        (a tail each); every 10 s, always. Sent together they pay \emph{one} promotion + \emph{one} tail.
  \item \textbf{Round trips before a byte}: DNS (1, unless cached) + TCP (1) + TLS 1.3 (1; TLS 1.2: 2) +
        request (1) — 4 RTTs. QUIC merges transport + TLS (3); a reused connection needs 1. Big
        bodies then pay \textbf{slow start}: the window doubles per RTT from $\approx$ 14 KB (RFC 6928).
  \item \textbf{Loss} on cellular is common: with HTTP/2 over TCP one lost packet stalls every stream;
        HTTP/3 stalls only its own stream and \textbf{migrates} across Wi-Fi $\leftrightarrow$ cellular.
\end{itemize}

\hd{Connection reuse — the biggest lever}
\begin{itemize}
  \item \textbf{One shared \ct{URLSession}} per configuration: its pool holds the warm
        connections. A session per request = full handshake every time.
  \item \textbf{HTTP/2}: one connection, many streams — no sharding. \textbf{Coalescing} (RFC 9113
        §9.1.1): a connection may serve another host if its certificate covers that host (browsers
        also want the same IP); the server refuses with \textbf{421 Misdirected Request}. Put
        \ct{api.} and \ct{img.} on one cert and one IP.
  \item \textbf{Idle pool timeout is not API}: Apple DTS — typically 30 s on HTTP/1.1, longer on
        HTTP/2, \textbf{a few seconds on WWAN}. ``Warm'' after a pause often means a new handshake.
\end{itemize}

\hd{Example — where did the time go?}
\begin{lstlisting}[language=SwiftSheet]
func urlSession(_ s: URLSession, task: URLSessionTask,
                didFinishCollecting m: URLSessionTaskMetrics) {
  for t in m.transactionMetrics where t.resourceFetchType == .networkLoad {
    func ms(_ a: Date?, _ b: Date?) -> Int {   // -1 = phase skipped
      guard let a, let b else { return -1 }
      return Int(b.timeIntervalSince(a) * 1000) }
    log("dns",  ms(t.domainLookupStartDate, t.domainLookupEndDate),
        "tcp",  ms(t.connectStartDate, t.secureConnectionStartDate),
        "tls",  ms(t.secureConnectionStartDate, t.secureConnectionEndDate),
        "ttfb", ms(t.requestStartDate, t.responseStartDate),
        "reused", t.isReusedConnection,        // iOS 10
        "cell", t.isCellular, t.networkProtocolName ?? "?") // iOS 13
  } }
\end{lstlisting}

\columnbreak

\hd{Fewer, earlier, bigger}
\begin{itemize}
  \item \textbf{Kill waterfalls}: 5 dependent calls = 5+ RTTs. Parallelise what is independent;
        a screen-level endpoint / BFF or GraphQL for what is not.
  \item \textbf{Prefetch} what the next screen almost surely needs, while the radio is already
        up; \textbf{batch} analytics and flush on a user-triggered request.
  \item \textbf{Priorities}: \ct{URLSessionTask.priority} (0–1, default 0.5) — a hint;
        \textbf{cancel} off-screen image loads; defer bulk work to a background session.
  \item \textbf{Warm up early}: start the first API call (its DNS + TLS) at launch, in parallel
        with UI setup.
\end{itemize}

\hd{Adapt to the network}
\ct{NWPathMonitor}: \ct{isExpensive} (cellular, hotspot), \ct{isConstrained} = \textbf{Low Data Mode}
(iOS 13) $\to$ skip prefetch, smaller images; per request \ct{allowsConstrainedNetworkAccess}.
\ct{waitsForConnectivity} (iOS 11) waits for a path instead of failing. \ct{timeoutIntervalForRequest}
(60 s) = max \emph{idle gap}; \ct{timeoutIntervalForResource} (7 days) = whole task.

\hd{Measure}
Metrics per phase (example) from the field · \textbf{Instruments $\to$ Network $\to$ HTTP Traffic}
(Xcode 13): sessions $\to$ tasks $\to$ transactions, on device · \textbf{Network Link Conditioner}
(Settings $\to$ Developer): 3G, LTE, \emph{Very Bad Network}, 100\% loss.

\hd{Traps}
\begin{itemize}
  \trap{``Fast on office Wi-Fi'' — test on Link Conditioner and a real cell.}
  \trap{``Need more bandwidth'' — API calls are latency-bound: count RTTs.}
  \trap{A new \ct{URLSession} per request — a handshake every time.}
  \trap{Polling every 10 s — the radio never leaves the tail. Push or batch.}
  \trap{\ct{timeoutIntervalForRequest} resets on every packet: not a total.}
\end{itemize}

\hd{Remember}
\textbf{``Wake once, round-trip less, reuse always, measure the phases.''}


\end{multicols}

\noindent{\footnotesize\color{sheetGrey}\textit{Related:} app-size-and-energy (tail diagram, batching,
isDiscretionary) · http-deep (h2 vs h3) · tls-pki (0-RTT) · resumable-uploads · nat-ipv4-ipv6}

\end{document}
